\documentclass[../../main.tex]{subfiles}
\graphicspath{{\subfix{../../image/}}}

\begin{document}

\section{独立同分布场合的极限定理}\label{section:ch05-独立同分布极限定理}

\subsection*{一、独立和问题}

在 \S1 中，我们讨论了 Bernoulli 试验场合事件 $A$ 出现次数 $\mu_n$ 的极限行为，曾指出 $\mu_n$ 可以表示为 $n$ 个独立随机变量之和（以后简称``独立和''），并对它证明了大数定律及中心极限定理，后来又看到这些定理有重要应用。这里自然会提出这样一个问题：这些性质是否只在 Bernoulli 试验场合才具有？

研究表明，许多独立和具有类似的性质，本节就要进一步讨论这个问题。

独立和的问题经常出现，例如测量一物体的某种尺寸，如测量一个圆柱体的直径 $d$，通常采用的办法是对它进行 $n$ 次测量，得到数值 $\xi_1,\xi_2,\cdots,\xi_n$，然后采用平均值
\begin{align*}
\eta_n=\frac{\xi_1+\xi_2+\cdots+\xi_n}{n}
\end{align*}
作为 $d$ 的数值。我们知道，测量时有各种随机因素影响，因此其结果带有随机性，这时 $\eta_n$ 是随机变量之和，如果各次测量是独立的，$\eta_n$ 便是独立和。为了说明上面所用办法的合理性就必须研究独立和。但这里的 $\xi_n$ 不服从 Bernoulli $0$--$1$ 分布，因此已不是 Bernoulli 试验场合的问题了。

在数理统计中已经把上述做法一般化。为了研究总体（它通常描述我们感兴趣的某一类现象）的某些特征，就对总体进行若干次观察以得到一批观察值 $\xi_1,\xi_2,\cdots,\xi_n$，并称它们是一个容量为 $n$ 的样本。再利用这个样本来构造各种统计量，例如
\begin{align*}
\frac{\xi_1+\xi_2+\cdots+\xi_n}{n}\quad\text{或}\quad \frac{\xi_1^{2}+\xi_2^{2}+\cdots+\xi_n^{2}}{n}
\end{align*}
以对总体的相应特征作各种推断。虽然每次观察得到的是具体的数值，但是为了比较各个统计量或各种推断方法的优劣，有必要把这些观察看作是某随机变量 $\xi$ 的观察值，通常假定 $\xi_1,\xi_2,\cdots,\xi_n$ 是相互独立的，且它们与 $\xi$ 具有相同的概率分布，这时上面的两个统计量便都是独立同分布（independent and identically distribution 简称 i.i.d.）的随机变量之和。

独立和的问题在许多实际问题中也出现，例如在计算电车整流站的电力负荷时，就遇到独立和问题，因为整流站的电力负荷等于各电车使用电力之和，每辆电车在某时刻的用电量是随机的，作为初步近似，可以假定各电车的用电量是相互独立的，因此这里遇到的正是独立和。在车间用电问题中，若有多类车床，用电量各不相同，则总用电量也需通过独立和来计算。

可见独立和问题经常遇到，而且各加项一般都不是 $0$--$1$ 分布。本节专门讨论各个加项服从相同分布的场合，这是实际工作（特别是在数理统计）中最常碰到的。从数学方面来看，这是最简单、最基本、最便于处理的，而且所用的处理方法可以相当方便地运用到更一般的场合。

在上章 \S5 中，我们讨论过独立的服从同类型分布的随机变量之和的分布问题，证明了某些分布的再生性，使用的工具是特征函数。现在我们要处理的问题与那里有很大区别，首先我们将对很一般的分布进行讨论，因而再生性通常都不满足；其次，不是对固定的 $n$ 进行讨论，而是讨论 $n\to\infty$ 时的情况，即研究极限定理。从数学的角度来看，它们可以看作是 Bernoulli 试验场合极限定理的推广，这里也是研究大数定律与中心极限定理。

所使用的工具还是特征函数。我们已经看到它很适合于处理独立和问题，有了上节的连续性定理，我们将进一步看到，它也很适合于处理极限分布问题。事实上，正是由于特征函数这一有力工具的使用，使得所有古典极限定理在短期内便得到了完满的解决。Laplace 就已经知道并应用了特征函数，俄国数学家 Lyapunov（1857—1918）最先发现并证明了收敛于正态分布的连续性定理（但并未明确叙述），从那时起，特征函数的理论不断得到完善。在这当中，法国数学家 Lévy 有突出的贡献。现在，特征函数法已经成了概率论的基本方法之一。

\subsection*{二、Khinchin 大数定律}

在 \S1 中，我们已经通过 Chebyshev 不等式建立起多种大数定律，那里都假定了方差的存在性，但是在独立同分布场合，并不需要有这个要求，这就是有名的 Khinchin（1894—1959）大数定律告诉我们的。用特征函数作为工具，这个定理很容易证明。

\begin{theorem}[Khinchin 大数定律]\label{theorem:ch05-Khinchin大数定律}
设 $\xi_1,\xi_2,\cdots,\xi_n,\cdots$ 是相互独立的随机变量序列，它们服从相同的分布，且具有有限的数学期望
\begin{align*}
a=\boldsymbol{E}\xi_n,
\end{align*}
则对任意的 $\varepsilon>0$，有
\begin{align}
\lim_{n\to\infty}P\!\left\{\left|\frac1n\sum_{i=1}^{n}\xi_i-a\right|<\varepsilon\right\}=1.
\end{align}
\end{theorem}

\begin{proof}
\end{proof}

显然，Bernoulli 大数定律是 Khinchin 大数定律的特殊情况。

Khinchin 大数定律在理论及应用中，特别是在数理统计中，十分重要，下面通过两个例子来略加说明。

\begin{example}[矩估计的相合性]\label{example:ch05-矩估计相合性}
假定总体 $\xi$ 的均值 $m_1$ 未知，通常的做法是对 $\xi$ 进行 $n$ 次独立重复观察，得到样本 $\xi_1,\xi_2,\cdots,\xi_n$，并以它们的平均值
\begin{align*}
A_1=\frac1n\sum_{i=1}^{n}\xi_i
\end{align*}
作为 $m_1$ 的估计量，这样做法的根据之一是依 Khinchin 大数定律应有
\begin{align*}
A_1\xrightarrow{P}m_1.
\end{align*}
它表明，当样本容量 $n$ 很大时，$A_1$ 作为 $m_1$ 的估计量是合理的。这个性质在数理统计中称为\textbf{相合性}，是选择估计量的最起码标准之一。

更为重要的是，根据 Khinchin 大数定律，若总体的 $k$ 阶原点矩 $m_k=\boldsymbol{E}\xi^{k}$ 存在，这时样本的 $k$ 阶原点矩
\begin{align*}
A_k=\frac1n\sum_{i=1}^{n}\xi_i^{k}
\end{align*}
作为 $m_k$ 的估计量也成立
\begin{align}
A_k\xrightarrow{P}m_k,
\end{align}
即样本 $k$ 阶原点矩 $A_k$ 是总体 $k$ 阶原点矩 $m_k$ 的相合估计量。利用中心矩与原点矩的简单关系，立即可证中心矩也有类似性质。因此 Khinchin 大数定律保证了矩估计的相合性。
\end{example}

\begin{example}[用 Monte Carlo 方法计算定积分]\label{example:ch05-MonteCarlo定积分}
为计算积分
\begin{align}
J=\int_{a}^{b}g(x)\,\mathrm{d}x
\end{align}
可以通过下面概率论方法实现。

任取一列相互独立的、都具有 $[a,b]$ 中均匀分布的随机变量 $\{\xi_i\}$，则 $\{g(\xi_i)\}$ 也是一列相互独立相同分布的随机变量，而且
\begin{align*}
\boldsymbol{E}g(\xi_i)=\frac{1}{b-a}\int_{a}^{b}g(x)\,\mathrm{d}x=\frac{J}{b-a}.
\end{align*}
既然
\begin{align}
J=(b-a)\cdot\boldsymbol{E}g(\xi_i),
\end{align}
因此只要能求得 $\boldsymbol{E}g(\xi_i)$，便能得到 $J$ 的数值。

为求 $\boldsymbol{E}g(\xi_i)$，自然想到大数定律，因为
\begin{align}
\frac{g(\xi_1)+g(\xi_2)+\cdots+g(\xi_n)}{n}\xrightarrow{P}\boldsymbol{E}g(\xi_i).
\end{align}
这样一来，只要能生成随机变量序列 $\{g(\xi_i)\}$ 就能对积分进行数值计算，而生成 $\{g(\xi_i)\}$ 的关键是要生成相互独立相同分布的 $\{\xi_i\}$，这里的 $\xi_i$ 均服从 $[a,b]$ 上的均匀分布。

现在已经可以把上述想法变成现实。这就是在电子计算机上产生服从均匀分布 $[a,b]$ 的随机数 $\{\xi_i\}$ 并利用上述两式估算 $J$，这种做法与我们在蒲丰投针问题中通过投针计算圆周率 $\pi$ 的做法是一致的。这种通过概率论的想法构造模型从而实现数值计算的方法，正如第一章 \S4 所言，已形成一种新的计算方法——\textbf{概率计算方法}，亦称 Monte Carlo 方法，它在原子物理、公用事业理论中发挥了不少作用，这个方法的理论根据之一就是大数定律。

至于计算积分，Monte Carlo 方法的实用场合是计算重积分
\begin{align}
I=\int_{K}g(P)\,\mathrm{d}P,
\end{align}
其中 $P$ 是 $m$ 维空间的点，当 $m$ 较大时，用 Monte Carlo 方法比一般数值法有优点，主要是它的误差与维数 $m$ 无关。
\end{example}

\subsection*{三、中心极限定理}

我们转而考虑如何把积分极限定理推广到相互独立相同分布，但分布函数为任意的随机变量序列的场合，这类问题在实际应用中非常重要。

若 $\xi_1,\xi_2,\cdots,\xi_n,\cdots$ 是一串相互独立相同分布的随机变量序列，且
\begin{align}
\boldsymbol{E}\xi_k=\mu,\qquad \boldsymbol{D}\xi_k=\sigma^{2}.
\end{align}
我们来讨论标准化随机变量和
\begin{align}
\zeta_n=\frac{1}{\sigma\sqrt{n}}\sum_{k=1}^{n}(\xi_k-\mu)
\end{align}
的极限分布。

Lindeberg 与 Lévy 建立了下列中心极限定理。

\begin{theorem}[Lindeberg--Lévy]\label{theorem:ch05-LindebergLevi}
对于标准化和，若 $0<\sigma^{2}<\infty$，则
\begin{align}
\lim_{n\to\infty}P\{\zeta_n<x\}=\frac{1}{\sqrt{2\pi}}\int_{-\infty}^{x}\mathrm{e}^{-t^{2}/2}\,\mathrm{d}t.
\end{align}
\end{theorem}

\begin{proof}
记 $\xi_k-\mu$ 的特征函数为 $g(t)$，则 $\zeta_n$ 的特征函数为
\begin{align*}
\left[g\!\left(\frac{t}{\sigma\sqrt{n}}\right)\right]^{n}.
\end{align*}
由于 $\boldsymbol{E}\xi_k=\mu$，$\boldsymbol{D}\xi_k=\sigma^{2}$，故 $g'(0)=0$，$g''(0)=-\sigma^{2}$。因此
\begin{align}
g(t)=1-\frac12\sigma^{2}t^{2}+o(t^{2}).
\end{align}
所以
\begin{align}
\left[g\!\left(\frac{t}{\sigma\sqrt{n}}\right)\right]^{n}
=\left[1-\frac{1}{2n}t^{2}+o\!\left(\frac{t^{2}}{n}\right)\right]^{n}\to\mathrm{e}^{-t^{2}/2}.
\end{align}
由于 $\mathrm{e}^{-t^{2}/2}$ 是连续函数，它对应的分布函数为 $N(0,1)$，因此由逆极限定理知
\begin{align*}
P\{\zeta_n<x\}\to\frac{1}{\sqrt{2\pi}}\int_{-\infty}^{x}\mathrm{e}^{-t^{2}/2}\,\mathrm{d}t.
\end{align*}
定理证毕。
\end{proof}

用这个定理立即可以推出 De Moivre--Laplace 积分极限定理。

Lindeberg--Lévy 定理有广泛应用。在实际工作中，只要 $n$ 足够大，便可以把独立同分布的随机变量之和当作是正态变量。这种做法在数理统计中用得尤其普遍。

\begin{example}[大样本理论]\label{example:ch05-大样本理论}
在数理统计中，为对总体 $\xi$ 的许多未知特征进行推断，通常的做法是抽取一个容量为 $n$ 的样本 $\xi_1,\xi_2,\cdots,\xi_n$，把它们看作独立同分布随机变量，为进一步提取信息，还构造一个或几个统计量 $g_k(\xi_1,\xi_2,\cdots,\xi_n)$，$k=1,2,\cdots,m$ 作为主要工具。在推断中通常需要知道这些统计量的分布，事实上却又十分难求，一种解决途径是借助于大样本理论，即在样本容量很大时，求这些统计量的渐近分布。这时多半利用 Lindeberg--Lévy 中心极限定理。

事实上，若 $m_{2k}=\boldsymbol{E}\xi_n^{2k}$ 存在，则由 Lindeberg--Lévy 中心极限定理，有
\begin{align*}
A_k=\frac1n\sum_{i=1}^{n}\xi_i^{k}\text{ 的分布渐近于 }N\!\left(m_k,\frac{m_{2k}-m_k^{2}}{n}\right).
\end{align*}
特别地，当 $k=1$ 时，若 $\boldsymbol{E}\xi_n=\mu$，$\boldsymbol{D}\xi_n=\sigma^{2}$ 存在，则
\begin{align*}
\frac1n\sum_{i=1}^{n}\xi_i\text{ 的分布渐近于 }N\!\left(\mu,\frac{\sigma^{2}}{n}\right).
\end{align*}
这些结论在统计学的几大分支：点估计、置信区间及假设检验中都经常用到。
\end{example}

\begin{example}[正态随机数的产生]\label{example:ch05-正态随机数产生}
在 Monte Carlo 方法中经常需要产生服从正态分布的随机数，但是一般计算机只备有产生 $[0,1]$ 均匀分布随机数（实际上是伪随机数）的程序。怎样通过 $[0,1]$ 均匀分布的随机数来产生正态随机数呢？这有多种途径，最常用的是利用上述定理来实现。

设 $\xi_1,\xi_2,\cdots,\xi_n,\cdots$ 是相互独立、均服从 $[0,1]$ 均匀分布的随机变量，这时\cref{theorem:ch05-LindebergLevi}的条件得到满足，故 $\xi_1+\xi_2+\cdots+\xi_n$ 渐近于正态变量。一般 $n$ 取不太大的值就可满足实际要求。图 5.3.1 中给出了 $n=1,2,3$ 时的图像。在 Monte Carlo 方法中，一般取 $n=12$，并用下式得到新的随机数序列：
\begin{align}
\eta_k=\sum_{i=1}^{12}\xi_{12(k-1)+i}-6,\qquad k=1,2,\cdots.
\end{align}
显然 $\{\eta_k\}$ 也是独立随机数序列，而且 $\boldsymbol{E}\eta_k=0$，$\boldsymbol{D}\eta_k=1$。经过检验证明，这时 $\eta_k$ 的渐近正态性已能满足一般精度要求。

\begin{figure}[H]
\centering
\includegraphics[width=0.6\textwidth]{ch05-fig-convolution.png}
\caption{均匀分布卷积}
\label{figure:ch05-均匀分布卷积}
\end{figure}
\end{example}

\begin{example}[近似数定点运算的误差分析]\label{example:ch05-误差分析}
数值计算时，任何数 $x$ 都只能用一定位数的有限小数 $y$ 来近似，这就产生了一个误差 $\xi=x-y$，在下面讨论中，我们假定参加运算的数都用十进制定点表示，每个数都用四舍五入的方法取到小数点后五位，这时相应的舍入误差可以看作是 $[-0.5\times10^{-5},0.5\times10^{-5})$ 上的均匀分布。

现在如果要求 $n$ 个数 $x_i$（$i=1,2,\cdots,n$）的和 $S$，在数值计算中就只能求出相应的有限位小数 $y_i$（$i=1,2,\cdots,n$）的和 $T$，并用 $T$ 作为 $S$ 的近似值。自然要问，这样做造成的误差 $\eta=S-T$ 是多少？

因为我们有
\begin{align*}
S=\sum_{i=1}^{n}x_i=\sum_{i=1}^{n}(y_i+\xi_i)=\sum_{i=1}^{n}y_i+\sum_{i=1}^{n}\xi_i,
\end{align*}
故
\begin{align*}
\eta=\sum_{i=1}^{n}\xi_i.
\end{align*}

一种传统的估计方法是这样的：由于
\begin{align*}
|\xi_i|\leqslant 0.5\times10^{-5},
\end{align*}
所以
\begin{align*}
|\eta|\leqslant\sum_{i=1}^{n}|\xi_i|\leqslant n\times0.5\times10^{-5}.
\end{align*}
以 $n=10\,000$ 为例，所得的误差估计为
\begin{align}
|\eta|\leqslant 0.05.
\end{align}

这种估计方法显然太保守，看来用概率论方法估计是适宜的。这时直接求 $\eta=\sum_{i=1}^{n}\xi_i$ 的分布不容易，但当 $n$ 较大时用极限定理作为工具，则能使问题很快得到解决。因为
\begin{align*}
\mu=\boldsymbol{E}\xi_i=0,\qquad \sigma=\sqrt{\boldsymbol{D}\xi_i}=\frac{0.5\times10^{-5}}{\sqrt{3}}.
\end{align*}
如果假定舍入误差 $\xi_i$ 是相互独立的，$n$ 又较大，那么用\cref{theorem:ch05-LindebergLevi}得到
\begin{align*}
P\!\left\{\left|\sum_{i=1}^{n}\xi_i\right|<k\sqrt{n}\sigma\right\}\approx\frac{1}{\sqrt{2\pi}}\int_{-k}^{k}\mathrm{e}^{-t^{2}/2}\,\mathrm{d}t.
\end{align*}
取 $k=3$ 时，上式右边为 $0.997$，因此我们能以 $99.7\%$ 的概率断言：
\begin{align}
|\eta|<3\times100\times\frac{0.5\times10^{-5}}{\sqrt{3}}=0.866\times10^{-3}.
\end{align}
这只是上限估计式中上限估计的 $60$ 分之一。

历史上，误差分析是概率论的重要生长点之一。19 世纪初德国数学家 Gauss 正是在研究测量误差时引进了正态分布并发展了有广泛应用的最小二乘法。至今这仍是概率论与生产实际有广泛联系的领域之一。
\end{example}

下面我们把中心极限定理推广到多变量的场合。

\begin{theorem}[多元中心极限定理]\label{theorem:ch05-多元中心极限定理}
若 $p$ 维随机向量 $\boldsymbol{\xi}_1,\boldsymbol{\xi}_2,\cdots,\boldsymbol{\xi}_n,\cdots$ 相互独立，具有相同的分布，其数学期望为 $\boldsymbol{\mu}$，协方差阵为 $\boldsymbol{\Sigma}$，则
\begin{align}
\boldsymbol{\eta}_n=\frac{(\boldsymbol{\xi}_1-\boldsymbol{\mu})+(\boldsymbol{\xi}_2-\boldsymbol{\mu})+\cdots+(\boldsymbol{\xi}_n-\boldsymbol{\mu})}{\sqrt{n}}
\end{align}
的极限分布为 $N(\boldsymbol{0},\boldsymbol{\Sigma})$。
\end{theorem}

\begin{proof}
对 $p$ 维列向量 $\boldsymbol{\lambda}$，构造
\begin{align}
\zeta_n=\frac{1}{\sqrt{n}}\sum_{i=1}^{n}\boldsymbol{\lambda}^{\mathrm{T}}(\boldsymbol{\xi}_i-\boldsymbol{\mu})=\boldsymbol{\lambda}^{\mathrm{T}}\boldsymbol{\eta}_n.
\end{align}
由于
\begin{align*}
\boldsymbol{E}\zeta_n=\frac{1}{\sqrt{n}}\sum_{i=1}^{n}\boldsymbol{\lambda}^{\mathrm{T}}(\boldsymbol{E}\boldsymbol{\xi}_i-\boldsymbol{\mu})=0,
\end{align*}
\begin{align*}
\boldsymbol{D}\zeta_n=\boldsymbol{E}\zeta_n^{2}
=\frac1n\boldsymbol{E}\!\left[\sum_{i=1}^{n}\boldsymbol{\lambda}^{\mathrm{T}}(\boldsymbol{\xi}_i-\boldsymbol{\mu})\cdot\sum_{j=1}^{n}(\boldsymbol{\xi}_j-\boldsymbol{\mu})^{\mathrm{T}}\boldsymbol{\lambda}\right]
=\frac1n\sum_{i=1}^{n}\boldsymbol{E}\!\left[\boldsymbol{\lambda}^{\mathrm{T}}(\boldsymbol{\xi}_i-\boldsymbol{\mu})(\boldsymbol{\xi}_i-\boldsymbol{\mu})^{\mathrm{T}}\boldsymbol{\lambda}\right]
=\frac1n\sum_{i=1}^{n}\boldsymbol{\lambda}^{\mathrm{T}}\boldsymbol{\Sigma}\boldsymbol{\lambda}
=\boldsymbol{\lambda}^{\mathrm{T}}\boldsymbol{\Sigma}\boldsymbol{\lambda},
\end{align*}
因此 $\zeta_n$ 是均值为 $0$，方差为 $\boldsymbol{\lambda}^{\mathrm{T}}\boldsymbol{\Sigma}\boldsymbol{\lambda}$ 的一维随机变量，由\cref{theorem:ch05-LindebergLevi}知它的分布函数收敛于 $N(0,\boldsymbol{\lambda}^{\mathrm{T}}\boldsymbol{\Sigma}\boldsymbol{\lambda})$。因此，若以 $f_n(t)$ 记 $\zeta_n$ 的特征函数，则由正极限定理知
\begin{align}
f_n(t)\to f(t,\boldsymbol{\lambda})=\exp\!\left(-\boldsymbol{\lambda}^{\mathrm{T}}\boldsymbol{\Sigma}\boldsymbol{\lambda}t^{2}/2\right)\qquad(n\to\infty).
\end{align}
而
\begin{align*}
f_n(t)=\boldsymbol{E}\mathrm{e}^{\mathrm{i}t\zeta_n}=\boldsymbol{E}\!\left\{\exp\!\left(\mathrm{i}t\boldsymbol{\lambda}^{\mathrm{T}}\boldsymbol{\eta}_n\right)\right\},
\end{align*}
因而
\begin{align*}
f_n(1)=\boldsymbol{E}\mathrm{e}^{\mathrm{i}\zeta_n}=\boldsymbol{E}\!\left\{\exp\!\left(\mathrm{i}\boldsymbol{\lambda}^{\mathrm{T}}\boldsymbol{\eta}_n\right)\right\}.
\end{align*}
它作为 $\boldsymbol{\lambda}$ 的函数，是 $\boldsymbol{\eta}_n$ 的特征函数。在上述极限式中，令 $t=1$，得到
\begin{align}
f_n(1)\to f(1,\boldsymbol{\lambda})=\exp\!\left(-\boldsymbol{\lambda}^{\mathrm{T}}\boldsymbol{\Sigma}\boldsymbol{\lambda}/2\right).
\end{align}
这正是 $p$ 维正态分布 $N(\boldsymbol{0},\boldsymbol{\Sigma})$ 的特征函数。因此由多维的连续性定理即得结论。
\end{proof}

\begin{example}\label{example:ch05-多项分布渐近正态}
服从 (3.2.6) 中多项分布的随机向量，可以看作 $n$ 个相互独立相同分布随机向量之和，由\cref{theorem:ch05-多元中心极限定理}可知多项分布渐近于正态分布，真正维数为 $r-1$。
\end{example}

\end{document}
